\documentclass[../../main.tex]{subfiles}
\begin{document}

\begin{center}
  {\fontsize{48pt}{58pt}\selectfont \textbf{1990}\normalsize\textbf{年度}}
\end{center}

\begin{table}[htb]
  \centering
  \caption{出題テーマと難易度}
  \begin{tabular}{|c|p{9cm}|c|}
    \hline
    \textbf{問題番号} & \textbf{テーマ} & \textbf{難易度} \\
    \hline
    第1問 & 和の極限と積分による評価 & \\
    \hline
    第2問 & 三次関数と区間外の値の比較 & \\
    \hline
    第3問 & 正八面体の切り口と正方形の穴 & \\
    \hline
    第4問 & 行列で定めた点列と三角形の面積 & \\
    \hline
    第5問 & 円に外接し楕円に内接する平行四辺形 & \\
    \hline
    第6問 & サイコロで作る小数と確率の極限 & \\
    \hline
  \end{tabular}
\end{table}

\begin{table}[htb]
  \centering
  \caption{解答}
  \renewcommand{\arraystretch}{1.4}
  \begin{tabular}{|c|c|p{9cm}|}
    \hline
    \textbf{問題番号} & \textbf{小問} & \textbf{答え} \\
    \hline
    第1問 & -- & $\displaystyle\lim_{n\to\infty}a_n=\infty$，$\displaystyle\lim_{n\to\infty}\dfrac{b_n}{a_n}=\dfrac{1}{\sqrt{2}}$ \\
    \hline
    第2問 & (1) & $h(x)=4x^3-3x$ \\
    \cline{2-3}
    & (2) & $|x|>1$で$|f(x)|<|h(x)|$（証明問題） \\
    \hline
    第3問 & (1) & 切り口の周の長さは$3$で一定（証明問題） \\
    \cline{2-3}
    & (2) & 通過させることができる \\
    \hline
    第4問 & (1) & \\
    \cline{2-3}
    & (2) & \\
    \hline
    第5問 & -- & 必要十分条件は$a^2+b^2=a^2b^2$ \\
    \hline
    第6問 & (1) & $\dfrac{8}{215}$ \\
    \cline{2-3}
    & (2) & $\dfrac{11}{30}\le\alpha\le\dfrac{37}{90}$ \\
    \hline
  \end{tabular}
\end{table}

\paragraph{試験形態}
6問．試験時間は要確認．

\end{document}
